\chapter{مقدمه}
هدف از فصل مقدمه، شرح مختصر مسئله تحقيق، اهميت و انگیزه محقق از پرداختن به آن موضوع به همراه اشاره‌اي كوتاه به روش و مراحل تحقيق است. مقدمه، اولين فصل از ساختار اصلی پايان نامه/ رساله بوده و زمينه اطلاعاتی لازم برای خواننده فراهم می آيد. در طول مقدمه بايد سعی شود موضوع تحقيق با زبانی روشن، ساده و بطور عميق و هدفمند به خواننده معرفی شود. این فصل بايد خواننده را مجذوب و اهميت موضوع تحقيق را آشکار سازد. در مقدمه بايد با ارائه سوابق، شواهد تحقيقی و اطلاعات موجود ( با ذکر منبع ) به روش منظم، منطقی و هدف‌دار، خواننده را جهت داد و به سوی راه حل مورد نظر هدايت کرد. مقدمه مناسب‌ترين جا برای ارائه اختصارات و بعضی توضيحات کلی است توضيحاتی که شايد نتوان در مباحث ديگر در مورد آنها توضيح داد.	
مقدمه، یکی از ارکان اساسی و اصلی پایان نامه است، مهمترین قسمت‌های آن عبارتند از : 

\section{عنوان تحقیق}
 باید شناختی دقیق و روشن از حوزه ی موضوع تحقیق را عرضه دارد و خالی از هرگونه ابهام و پیچیدگی باشد. 

\section{مسئله تحقیق}
وظیفه اصلی مقدمه بیان این مطلب به خواننده است که چرا انجام تحقیق را به عهده گرفته‌اید. اگر دلیل شما برای انجام این کار پاسخگویی به سؤال مورد علاقه‌تان است، با مشکل زیادی روبه‌رو نخواهید بود. یکی از بهترین روش‌ها برای نوشتن مقدمۀ‌ يك پایان‌نامه، طرح پرسش یا پرسش‌هایی مهم و اساسی است که کار تحقیقاتی شما از آغاز تا پایان قصد پاسخ دادن به آن را دارد. گاهی می‌توانید ابتدا اهمیت موضوع را بیان و سپس پرسش خود را در آن موضوع مطرح کنید.

\section{تاریخچه ای از موضوع تحقیق}
به طور کلی تشریح روند های تحقیقاتی در محدوده مورد مطالعه، مستلزم ارجاع به کارهای دیگران است. بعضی از نویسندگان برای کارهای دیگران هیچ اعتباری قائل نمی‌شوند و در مقابل، بعضی دیگر از نویسندگان در توصیف کارهای دیگران، بسیار زیاده‌روی می‌کنند. اکثر مواقع ارجاع به مقالات دو سال قبل از کارتان، بهتر از نوشتن سطرها مرجع است. در این قسمت مختصری از نظرات و تحقیقات مربوط به موضوع و یا مسایل و مشکلات حل نشده در این حوزه و همچنین توجه و علاقه جامعه به این موضوع اشاره می‌شود. 

\section{تعریف موضوع تحقیق}
در این قسمت محقق موضوع مورد علاقه و یا نیاز احساس شده خود را در حوزه تحقیق بیان می‌دارد و عوامل موجود در موقعیت را تعریف و تعیین می‌کند

\section{هدف یا هدف‌های کلی و آرمانی تحقیق}
این قسمت باید با جملات مثبت و کلی طرح شود و از طولانی شدن مطالب پرهیز شود 

\section{روش انجام تحقیق}
در این قسمت پژوهشگر روش کاری خود را بیان می‌دارد و شیوه‌های گوناگونی را که در گردآوری مطالب خود به کار برده ذکر می‌کند و همچنین اگر روش آماری خاصی را در تهیه و تدوین اطلاعات به کار برده آن شیوه را بیان می‌کند.

\section{نوآوری، اهمیت و ارزش تحقیق}
در این قسمت نوآوری علمی و عملی تحقیق که محقق به آن دست خواهد یافت، بحث می‌شود.  ممکن است لازم باشد تا از برخی نمودار ها در این بخش استفاده شود. توجه شد، استفاده از شکل در این بخش تنها و تنها به جهت ارائه مثالی برای استفاده از نمودار ها و شکل ها می باشد. طبیعتاً به صلاح‌دید نگارنده شکل‌ها و نمودار‌ها می توانند در بخش های مختلف مورد استفاده قرار گیرند.